\documentclass{sp}

% The \pdf* commands provide metadata for the PDF output.
% Do not use LaTeX style / commands like \emph{} inside these.
\pdfauthor{Author Full Name(s)}
\pdftitle{Full title}
\pdfkeywords{XXX, XXX}

% Optional short title inside square brackets, for the running headers.
% If no short title is given, no title appears in the headers.
%\title[Adjective ordering preferences]{On the grammatical source of adjective ordering preferences%
\title[Practical introduction to RSA]{A practical introduction to the Rational Speech Act modeling framework%
	\thanks{We thank XXX.}}

% Optional short author inside square brackets, for the running headers.
% If no short author is given, no authors print in the headers.
\author[]{% As many authors as you like, each separated by \AND. %%% uncomment to de-anonymize
%	\spauthor{author1 \\ \institute{affiliation1}} \AND
%	\spauthor{author2 \\ \institute{affiliation2}} \AND
%	\spauthor{author3 \\ \institute{affiliation3}}
}

\usepackage{linguex}
\usepackage{qtree}
\qtreecenterfalse
\usepackage{tree-dvips}
\usepackage{phonetic}
\usepackage{xcolor}
\usepackage{pifont}
\usepackage{lineno}
\usepackage{graphicx}
%\usepackage{qdotbranch}
\usepackage{booktabs}
\usepackage{multirow}
\usepackage{CJK}
\usepackage{tikz}
\usepackage{amsfonts} 
\usepackage{amsmath}
\usepackage{hyperref}

\def\url#1{\expandafter\string\csname #1\endcsname}

\newcommand{\gcs}[1]{\textcolor{blue}{[gcs: #1]}}  
\newcommand{\mf}[1]{\textcolor{orange}{[mf: #1]}}  

\newcommand{\type}[1]{\ensuremath{\left \langle #1 \right \rangle }}
\newcommand{\lam}{\ensuremath{\lambda}}
\newcommand{\sem}[1]{\ensuremath{[\![#1]\!]}}

\renewcommand{\firstrefdash}{}

\begin{document}

\maketitle

\begin{abstract}
	Recent advances in computational cognitive science (i.e., simulation-based probabilistic programs) have paved the way for significant progress in formal, implementable models of pragmatics. Rather than describing a pragmatic reasoning process, these models articulate and implement one, deriving both qualitative and quantitative predictions of human behavior---predictions that consistently prove correct, demonstrating the viability and value of the framework. The current paper provides a practical introduction to the Bayesian Rational Speech Act modeling framework, serving as a companion piece to the hands-on web-book at \href{https://www.problang.org}{www.problang.org}. After providing a conceptual overview of the framework, we then walk readers through the resources of the web-book.
\end{abstract}

\begin{keywords}
	XXX, XXX
\end{keywords}

\tableofcontents

\section{Introduction}

Much work in formal, compositional semantics follows the tradition of positing systematic but inflexible theories of meaning. However, in practice, the meanings we derive from language are heavily dependent on nearly all aspects of context, both linguistic and situational. To formally explain these nuanced aspects of meaning and better understand the compositional mechanism that delivers them, recent work in formal pragmatics recognizes semantics not as one of the final steps in meaning calculation, but rather as one of the first. Within the Bayesian Rational Speech Act (RSA) framework \citep{goodmanfrank2016,frankejaeger2016}, speakers and listeners reason about each other's reasoning about the literal interpretation of utterances. The resulting interpretation necessarily depends on the literal interpretation of an utterance, but is not necessarily wholly determined by it. This move---reasoning about likely interpretations---provides ready explanations for complex phenomena ranging from metaphor \citep{kaoetal2014metaphor} and hyperbole \citep{kaoetal2014} to the specification of thresholds in degree semantics \citep{lassitergoodman2013}.

The probabilistic pragmatics approach leverages the tools of structured probabilistic models formalized in a stochastic $\lambda$-calculus to develop and refine a general theory of communication. The framework synthesizes the knowledge and approaches from diverse areas---formal semantics, Bayesian models of inference, formal theories of measurement, philosophy of language, etc.---into an articulated theory of language in practice. These new tools yield broader empirical coverage and richer explanations for linguistic phenomena through the recognition of language as a means of communication, not merely a vacuum-sealed formal system. By subjecting the heretofore off-limits land of pragmatics to articulated formal models, the rapidly growing body of research both informs pragmatic phenomena and enriches theories of linguistic meaning.

These models are particularly well-suited for capturing XXX Goodies list of RSA

The current paper offers a practical introduction to the modeling framework, serving as a companion piece to the hands-on web-book at \href{https://www.problang.org}{www.problang.org}. We begin in Section \ref{overview} with a high-level overview of RSA, walking through its basic implementation and the philosophical foundations that informed the architectural choices. We then explore variations to the basic architecture in Section \ref{variations}, surveying technological innovations that have allowed for broader empirical coverage. Section \ref{practicalities} discusses common practical considerations faced by the working modeler. In Section \ref{limitations}, we explore limitations of the current framework, which also serve as guidance for future extensions. Section \ref{summary} concludes.


\section{High-level overview of RSA} \label{overview}

The RSA framework views language understanding as a process of recursive social reasoning between speakers and listeners: listeners interpret the utterances they hear by reasoning about how speakers generate them; speakers choose their utterances by reasoning about how listeners interpret them. In the basic, vanilla RSA model from \cite{frankgoodman2012}, this recursion involves three layers of inference. Typically formulated as statements of conditional probability, as in (\ref{L0}), (\ref{S1}), and (\ref{L1}), these inference layers correspond to models of speakers and listeners.

\begin{equation} \label{L0}
P_{L_0}(s \mid u) \propto \sem{u}(s)
\end{equation}
\begin{equation} \label{U}
U_{S_1}(u; s) = \textrm{log}P_{L_0}(s \mid u) - C(u)
\end{equation}
\begin{equation} \label{S1}
P_{S_1}(u \mid s) \propto \textrm{exp}(\alpha \cdot U_{S_1}(u;s))
\end{equation}
\begin{equation} \label{L1}
P_{L_1}(s \mid u) \propto P_{S_1}(u \mid s) \cdot P(s)
\end{equation}

The reasoning grounds out in the naive, literal listener, $L_0$, who interprets utterances
according to their literal semantics. The semantics of utterance $u$ is captured in meaning
function $\sem{u} \colon s \mapsto [0;1]$ which maps states to truth values. We assume binary
truth-values here, but the formalism works just as well for non-binary, e.g., fuzzy semantics
\mf{insert refs: Degen on overspecification, Steinert on vagueness, ...?}. According to
\eqref{L0}, $L_0$ hears some utterance $u$ and infers the state of the world $s$ that $u$ was
meant to describe by restricting the set of possible states to just those that are compatible
with the literal, truth-functional semantics of $u$, returning a uniform probability
distribution over the states $s$ that $u$ maps to \texttt{true}.

One layer up, a pragmatic speaker, $S_1$, chooses utterances in proportion to their utility $U_{S_{1}}$. Utterances are useful to the extent that they maximize the probability that $L_0$ will infer the correct $s$ on the basis of $u$, while minimizing the cost of $u$ ($C(u)$; speakers aim to be efficient). So, when selecting utterances, $S_1$ considers their effect on interpretation (i.e., on $L_0$'s resulting beliefs); utterances that are most likely to lead $L_0$ to the correct belief are most likely to be chosen by $S_1$.

At the top layer of inference, the pragmatic listener, $L_1$, interprets utterances to infer the true state of the world. However, unlike $L_0$, who reasons directly about the utterance semantics, $L_1$ reasons instead about the process that generated the utterance; that process is the speaker $S_1$. $L_1$ thus infers $s$ on the basis of $u$ by reasoning about the probability that $S_1$ would have chosen $u$ to signal $s$ to $L_0$; the higher that probability, the more likely $L_1$ is to conclude that $S_1$ indeed intended to communicate $s$. Because $L_1$ reasons about $S_1$, who in turn reasons about the literal semantics in $L_0$, $L_1$'s interpretation is affected by the semantics of $u$, albeit only indirectly via the $S_1$ layer. This space between the semantics (i.e., $L_0$) and the resulting interpretation (i.e., the posterior beliefs of $L_1$) is where pragmatics enters. To see how, it would help to consider a concrete communication scenario.

In its initial formulation, \cite{frankgoodman2012} used the RSA framework to model referent choices in efficient communication. Suppose we are in a world as in Figure \ref{ref-game} with three objects: a blue square, a blue circle, and a green square. Suppose further that a speaker is trying to signal a single object in this world to a cooperative listener, and that the speaker can only use a single-word utterance to do so. The utterances available to the speaker include ``blue'', ``green'', ``square'', and ``circle''; the possible states the listener might infer correspond to the three objects: blue-square, blue-circle, and green-square. We have the expected truth-functional semantics for the utterances: ``blue'' maps blue-square and blue-circle to true but green-square to false, ``green'' maps blue-square and blue-circle to false but green-square to true, etc.

\begin{figure}[t]
\centering
\includegraphics[width=3in]{rsa-scene.pdf}
\caption{An example referential-communication game from \cite{frankgoodman2012}.}
\label{ref-game}
\end{figure}

With the semantics as stated, $L_0$ interprets utterances to return uniform probability distributions over the compatible states. Hearing the utterance ``blue'', $L_0$ returns a belief distribution over states that divides probability equally between blue-square and blue-circle, the only objects compatible with the semantics of ``blue''. Hearing ``circle'', $L_0$ returns a distribution with 100\% of the probability on blue-circle---$L_0$ is certain that the speaker intends to signal blue-circle. With $L_0$ formulated in this way, we have an agent who interprets utterances naively, according to the literal semantics. 

$S_1$ chooses utterances by simulating their effect on $L_0$. Suppose the speaker wants to communicate blue-circle to $L_0$. Two utterances, ``green'' and ``square'', stand no chance of communicating the intended state to $L_0$ and so they are ruled out. The other two utterances, ``blue'' and ``circle'', are both literally compatible with blue-circle, but one of the utterances is more likely to lead $L_0$ to the correct belief state. If the speaker were to utter ``blue'', we saw that $L_0$'s belief distribution would be evenly split between blue-square and blue-circle. In other words, ``blue'' has a 50\% chance of leading $L_0$ to the correct state. On the other hand, ``circle'' has a 100\% chance of leading $L_0$ to the correct state. Assuming equal utterance costs, ``circle'' is thus twice as useful to the speaker as ``blue'' for communicating blue-circle to $L_0$, and $S_1$ reflects this asymmetry in utility by assigning twice as much probability to ``circle'' in the posterior distribution over utterance choices.\footnote{With no scaling. XXX} For communicating the blue-square state, both ``blue'' and ``square'' have a 50\% chance of leading $L_0$ to the correct state, so both utterances are equally likely and thus equally probable in the $S_1$ posterior.

Within this simple reference-game scenario, $L_1$ reasons pragmatically about $S_1$ to break the symmetry in the semantics. Hearing ``blue'', $L_1$ will invert the $S_1$ model to determine which state (i.e., which object) the speaker is most likely trying to communicate. Had the speaker wanted to communicate the blue-circle state, we saw above that the speaker would have been more likely to utter ``circle''. But the speaker did not utter ``circle''; she uttered ``blue'' instead. This counterfactual reasoning leads $L_1$ to down-weight the probability of the state that ``circle'' could have uniquely picked out, since the speaker could have said ``circle'' but chose not to. As a result, more probability gets assigned to the blue-square state, the other state compatible with the semantics of the utterance. In this way, we capture the Gricean specificity implicature associated with uttering ``blue'' in a scenario as in Figure \ref{ref-game}: the speaker probably intends the blue square because if she wanted to communicate the blue circle she could have said ``circle''.

A key component of the RSA framework is the updating of prior beliefs: listeners use the utterances they hear as signal with which to update their beliefs about the world. In walking through the reference-game example above, we assumed a uniform prior over world states: before hearing the speaker's utterance, $L_1$ had no reason to suspect that any object was more or less likely to get referenced. With a different, more informative prior, it is possible to shift the qualitative predictions of $L_1$. Suppose that $L_1$ has reason to suspect that the blue-circle state is most likely to get referenced. Now there are two opposing pressures operating on $L_1$'s interpretation of the utterance ``blue" (i.e., on $L_1$'s posterior beliefs): the prior belief favoring blue-circle over blue-square and the specificity implicature discussed above that favors blue-square over blue-circle (if the speaker had wanted to reference blue-circle, she could have said ``circle''). With a sufficiently strong prior in favor of blue-circle, it is possible to override the specificity implicature so that $L_1$ interprets ``blue'' as confirming his suspicions and referring to blue-circle.

XXX scalar implicature?


\paragraph{Relation to Gricean maxims.} The vanilla RSA model, as described above in Equations
\eqref{L0}--\eqref{L1}, is a direct formalization of Grice's idea that listeners can infer
pragmatic meaning based on the assumption that speakers adhere to certain rules of behavior,
the so-called ``Maxims of Conversation'' \citep{Grice1975:Logic-and-Conve}. This relationship to Grice's maxims becomes
transparent if we rewrite the speaker's utterance-choice probability, starting from the
definition in \eqref{S1}.
%
\begin{align} \label{eq:S1-rewrite}
  P_{S_1(u\mid s)} & \propto \exp \left [ \alpha \left (\log P_{L_0}(s \mid u) - C(u) \right)  \right ] & \text{\textcolor{gray}{[definitions \eqref{U} \& \eqref{S1}]}} \\
  % & = \left [\exp \left ( \log P_{L_0}(s \mid u) - C(u) \right ) \right ]^\alpha & \text{\textcolor{gray}{[rules exponential function]}} \nonumber \\  
  & = \left (P_{L_0}(s \mid u) \  \frac{1}{\log C(u)} \right)^{\alpha} & \text{\textcolor{gray}{[rules exponential \& log]}} \nonumber \\
  & = \left ( \frac{\sem{u}(s)}{| \sem{u} |} \ \frac{1}{\log C(u)} \right )^{\alpha} & \text{\textcolor{gray}{[definition \eqref{L0}]}} \nonumber 
\end{align}
%
The reformulation in~\eqref{eq:S1-rewrite} shows how the speaker's choice probabilities are a
product of three factors, each corresponding to a central postulate concerning how cooperative and, arguably, rational speakers should tailor their contributions
to a conversation. Suppose we recast the components as follows:
$\text{Truth}(u,s) = \sem{u}(s)$, $\text{Informativity}(u) = |\sem{u}|^{-1}$, and
$\text{Economy(u)} = \log C(u)$. We can then rewrite the speaker rule as:\footnote{For
binary truth-values, $\alpha$ can be dropped from the factor $\text{Truth}$.}
\begin{align} \label{eq:S1-three-factor-formulation}
  P_{S_1(u\mid s)}   & \propto \text{Truth}(u,s) \ \text{Informativity}(u)^{\alpha} \ \text{Economy}(u)^{\alpha} % \propto \underbrace{\sem{u}(s)^{\alpha}}_{\text{Quantity}} \ \underbrace{| \sem{u} |^{-\alpha}}_{\text{Quantity}} \  \underbrace{(\log C(u))^{-\alpha}}_{\text{Manner}} & \text{\textcolor{gray}{[definition \eqref{L0}]}} \nonumber
\end{align}
These three factors directly capture (a formalization of) the Gricean maxims of Quality,
Quantity, and Manner. In effect, the RSA speaker assumes that speakers (i) maximize truth
(i.e., they utter no falsehoods if they can select at least one
true message), (ii) maximize informativity (all else equal) and (iii) maximize utterance
economy (all else equal). The more rational a speaker (i.e., the higher $\alpha$), the more
pronounced the latter two optimization tendencies become. This behavior reflects the special
status that Grice bestowed on the Maxim of Quantity: (binary) truthfulness is not affected by
the speaker's degree of optimizing informativity and costs.

The formulation in \eqref{eq:S1-rewrite} highlights the ease with which it is possible to define nearby
alternative models of speaker production, going beyond Grice's seminal work by
capturing additional psychological factors (e.g., from a more algorithmic perspective on how speakers
generate utterances). For instance, some applications might want to capture differences between
utterance alternatives that are not related to production economy or cost, but rather to the
ease with which an alternative utterance comes to mind. The latter factor of \emph{differential
  salience} is arguably not subject to optimization: it is not the case that the more rational
a speaker is, the more she would tend to select only the utterances which come to mind most
easily. We can implement differential salience (ease of retrieval) for instance in a
formulation like this:
%
\begin{align}
  \label{eq:S1-utterance-salience}
  P_{S_1(u\mid s)}^{\text{Salience}}   & \propto \text{Truth}(u,s) \ \text{Informativity}(u)^{\alpha} \ \text{Salience}(u)
\end{align}
An implementation of the speaker model in \eqref{eq:S1-utterance-salience} would implement the
factor $\text{Salience}(u)$ as an \emph{utterance prior}, suggesting an algorithmic
picture of utterance choice: the speaker searches for utterances to choose, based on a gradient
of how easy these utterances come to mind, then compares the available options (weighted by
their salience/prominence) based on the other factors relevant for communication: truth and
informativity. \mf{I'm not sure if this paragraph needs to go here; it might have a better
  place elsewhere; but stressing the 'algorithmic picture' somewhere might be good}

\paragraph{Informativity from alignment of beliefs.} We saw above that the vanilla RSA
model's definition of speaker's choice probabilities can be interpreted as a product of three
factors that correspond closely to Gricean maxims of Quality, Quantity and Manner. 
This is just one formal implementation; others are readily conceivable. Given the particular
importance of informativity for many pragmatic explanations, one would be justified in further dissecting
the way informativity is treated in the RSA model. In the following, we show that \eqref{U} \&
\eqref{S1} follows from a simple idea, which can be spelled out in a natural way: the speaker has some
belief about the relevant world states that gets captured in probability distribution $P_{S\text{-}Bel}$;
after an utterance of $u$, the literal listener's beliefs about the relevant world states are
represented as another probability distribution $P_{LL\text{-}Bel}(u)$. An utterance $u$ is better/more useful
than another utterance $u'$ if the listener's beliefs $P_{LL\text{-}Bel}(u)$ after hearing $u$ are
more aligned with the speaker's beliefs $P_{S\text{-}Bel}$ than the listener's beliefs
$P_{LL\text{-}Bel}(u')$ after hearing $u'$. If we spell out ``being more aligned with'' in terms of an
information-theoretic notion of divergence between probability distributions, the vanilla RSA
model is recovered. Behind it all, however, is simply the idea that the speaker prefers
utterances that best align speaker's and listener's relevant beliefs about the world.

A useful notion of ``alignment of probability distributions'' is the information-theoretic
measure of Kullback-Leibler divergence. KL-divergence is not symmetric, but considers one of
two to-be-compared probability distributions as the ground-truth or the objective function to
to be approximated. This assumption makes sense in our application to language because it is the speaker's
beliefs to which the listener should align. KL-divergence then measures divergence in terms of,
essentially, the expected number of extra bits needed to encode a signal that was generated
from the true distribution with the approximate distribution. \mf{leave out that last
  sentence?} The Kullback-Leibler divergence between (baseline/true) probability distribution
$P_{S\text{-}Bel}$ and (approximate/to-be-optimized) probability distribution
$P_{LL\text{-}Bel}(u)$ is defined as:
\begin{align}
  \label{eq:KL-divergence}
  \text{KL}(P_{S\text{-}Bel} \mid \mid P_{LL\text{-}Bel}(u)) = - \sum_{s} P_{S\text{-}Bel}(s) \ \log \frac{P_{S\text{-}Bel}(s)}{P_{LL\text{-}Bel}(u,s)}
\end{align}

Using KL-divergence, we can then state a more general definition of utterance utilities, to
replace \eqref{U}:
\begin{align}
  \label{eq:Utils-KL-based}
  U_{S_1}(u; s) = \text{KL}(P_{S\text{-}Bel} \mid \mid P_{LL\text{-}Bel}(u)) - C(u)
\end{align}
The formulation in \eqref{U} is a special case of the more general \eqref{eq:Utils-KL-based}
which follows if the speaker's beliefs about the relevant world states are degenerate, that is,
whenever the speaker knows the true world state $s^{*}$, so that $P_{S\text{-}Bel}(s^{*})=1$.
In that case, we have:
\begin{align*}
  \text{KL}(P_{S\text{-}Bel} \mid \mid P_{LL\text{-}Bel}(u)) & = - \sum_{s} P_{S\text{-}Bel}(s) \ \log \frac{P_{S\text{-}Bel}(s)}{P_{LL\text{-}Bel}(u,s)} \\
  & =  - \log\frac{1}{P_{LL\text{-}Bel}(u,s^*)} = \log P_{LL\text{-}Bel}(u,s^*)
\end{align*}


\section{Variations on vanilla} \label{variations}

We have seen how vanilla RSA can model pragmatic reasoning in simple reference games, as well as specificity implicatures more generally. In this section, we explore recent extensions of RSA meant to handle more complex language phenomena. We organize our discussion around deviations from vanilla RSA, highlighting novel technology and the phenomena it captures. As a point of reference for vanilla RSA, we repeat the $S_1$ utterance choice rule below, where a pragmatic speaker selects utterances in proportion to their utility in conveying some observed state of affairs to a naive listener who interprets the utterance according to its literal semantics---all while minimizing utterance cost. As we shall see, many of the extensions of RSA can be characterized in terms of their deviation from this simple utterance-choice rule.
\begin{equation} \label{S1-repeat}
P_{S_1}(u|s) \propto \textrm{exp}(\alpha (\textrm{log}P_{L_0}(s|\sem{u}) - C(u)))
\end{equation}

\subsection{Meaning inference}

Vanilla RSA assumes a fixed utterance semantics that is shared by both speakers and listeners. But what happens when we parameterize the interpretation function such that the utterance semantics is subject to change? This move is illustrated in \eqref{S1-lexical-uncertainty}, where the variable \textbf{x} determines the truth-functional mapping for an utterance $u$. Once named, this interpretation-fixing variable \textbf{x} becomes subject to active pragmatic reasoning. In other words, we can allow for uncertainty around \textbf{x}, such that the precise semantics of $u$ might be vague or underspecified in context. What results is the class of lexical uncertainty RSA models.
\begin{equation} \label{S1-lexical-uncertainty}
P_{S_1}(u|s, \textbf{x}) \propto \textrm{exp}(\alpha (\textrm{log}P_{L_0}(s|\sem{u}^{\textbf{x}}) - C(u)))
\end{equation}

\cite{lassitergoodman2013} use this technology to model the interpretation of vague gradable adjectives, whose meaning depends crucially on properties of the context, both linguistic and extra-linguistic. Take the adjective \emph{heavy}, which is true of some state (i.e., a weight) just in case the weight exceeds the relevant threshold for heaviness. However, we conclude that a \emph{heavy elephant} weighs substantially more than a \emph{heavy backpack}; moreover, the weight we attribute to the backpack is likely to vary depending on whether the speaker is a preschooler or an olympian. Thus, the threshold above which some state counts as heavy is unlikely to be a fixed, context-invariant value. 

\citeauthor{lassitergoodman2013} handle the uncertainty around the threshold value by parameterizing the meaning of utterances containing gradable adjectives: \emph{heavy} will be true of some state just in case the weight exceeds the relevant threshold (as before), where this threshold (\textbf{x} in \eqref{S1-lexical-uncertainty}) gets fixed by pragmatic reasoning:

\ex. \label{heavy-sem}
\sem{heavy}$^{\textbf{x}}$ = \lam s. weight(s) $>$ \textbf{x}

It is up to the pragmatic listener, $L_1$, to resolve the uncertainty surrounding \textbf{x}. To do so, $L_1$ hears the utterance and jointly infers both the state of the world (i.e., the weight of the state being described) and the relevant heaviness threshold. $L_1$ performs this inference just as in the vanilla model: which state and threshold would have been most likely to lead $S_1$ to select the utterance that $L_1$ encountered? In other words, $L_1$ simulates $S_1$'s behavior for the possible combinations of states and thresholds, then weights states and thresholds in proportion to the probability that they would lead $S_1$ to choose the utterance that was encountered.

This technology---parameterization of the utterance semantics subject to pragmatic reasoning---allows for a recasting of the division of labor between semantics and pragmatics. With lexical uncertainty, the semantic content of an utterance can be (at least partially) determined via pragmatic inference. \cite{bergenetal2016} seize on this innovation to model M-implicatures \citep{horn1984} and certain embedded implicatures \citep{hurford1974,chierchiaetal2012}. Rather than uncertainty around the semantics of a single utterance (e.g., a gradable adjective, as in the example above), \citeauthor{bergenetal2016} assume uncertainty around entire lexica. \gcs{need to say more about lexical uncertainty?}

Meaning inference also serves as a useful tool for modeling ambiguity resolution. \cite{scontrasgoodman2017} use an utterance-choice rule as in \eqref{S1-lexical-uncertainty} to model the resolution of distributive-collective ambiguities in plural predication. If a listener hears that \emph{the boxes are heavy}, the listener must jointly infer both the intended interpretation of the utterance (i.e., distributive, commenting on individual box weight, or collective, commenting on the total weight of the set) and the state of the world (i.e., the weights of the boxes). Returning to \eqref{S1-lexical-uncertainty}, \citeauthor{scontrasgoodman2017} treat \textbf{x} as the interpretation-resolving variable that determines whether $u$ receives a distributive or collective interpretation. \cite{savinellietal2017,savinellietal2018} use this technology to model the resolution of scope ambiguity, as in utterances like \emph{every horse didn't jump over the fence}. For \citeauthor{savinellietal2017}, \textbf{x} in \eqref{S1-lexical-uncertainty} serves to determine the relative scope of quantification elements at LF. 

Both of these applications---plural predication and quantifier scope ambiguity---highlight how meaning inference in RSA can serve as a useful shorthand for providing computational-level descriptions of the process of ambiguity resolution. However, while these models offer a promising hypothesis for how listeners reason pragmatically to resolve ambiguity, it bears noting that neither application attempts to model the compositional processes that give rise to the relevant ambiguities in the first place. We return to this issue in Section \ref{limitations}.


\subsection{QUD inference}

The technology we have considered up to now allows for the pragmatic enrichment of utterances: from ``blue'' to ``blue square'' or from vague ``expensive'' to ``expensive'' with a contextually-determined price threshold. However, given that $S_1$'s utility relies on successfully communicating the observed $s$ to $L_0$ (i.e., on informativity), our enrichments can only ever be literal---we have no way for ``blue'' to be interpreted as ``green''. However, many instances of everyday language use are, strictly speaking, non-literal. For example, if you hear that someone paid ``a million dollars'' for their coffee at the hipster hangout around the corner, you are unlikely to infer that the actual price paid was \$1,000,000. Rather, you infer that the price was high and that the speaker is, to put it mildly, less than thrilled.

To allow for enrichments beyond the literal meaning, we must broaden $S_1$'s goals beyond informativity with respect to $s$. One way of conceiving of speaker goals is in terms of the Question Under Discussion (QUD) the speaker aims to answer. Under certain theories of pragmatics, all discourse proceeds with respect to some QUD \citep{roberts2012}; utterances in the discourse must at least partially answer the QUD in order to be cooperative and felicitous. By allowing for a broader range of QUDs and uncertainty around which QUD is intended, vanilla RSA may be extended to capture uses of non-literal language. 

As was the case with meaning inference in the previous subsection, we can illustrate this innovation by highlighting its effect on the vanilla utterance-choice rule. In \eqref{S1-QUD}, $S_1$'s utility continues to break down into informativity and economy components. However, now \textbf{x} determines what it is that $S_1$ endeavors to communicate to $L_0$. Viewed as the QUD, \textbf{x} serves to map $s$ to the answer to \textbf{x} that $S_1$ would like to communicate to $L_0$.
\begin{equation} \label{S1-QUD}
P_{S_1}(u|s, \textbf{x}) \propto \textrm{exp}(\alpha (\textrm{log}P_{L_0}(F(s,\textbf{x})|\sem{u}) - C(u)))
\end{equation}

\cite{kaoetal2014} use this technology to model hyperbole, as in the coffee-price example above. To see how, consider some possible states of the world about which a speaker might want to communicate:

\ex. \label{hyperbole-states}
\emph{Possible states of the world for the coffee-price scenario}:\\
$s_1$: \{\texttt{affect}: positive, \texttt{price}: \$3\}\\
$s_2$: \{\texttt{affect}: positive, \texttt{price}: \$7\}\\
$s_3$: \{\texttt{affect}: positive, \texttt{price}: \$1,000,000\}\\
$s_4$: \{\texttt{affect}: negative, \texttt{price}: \$3\}\\
$s_5$: \{\texttt{affect}: negative, \texttt{price}: \$7\}\\
$s_6$: \{\texttt{affect}: negative, \texttt{price}: \$1,000,000\}

Notice that the world states in \ref{hyperbole-states} have two properties: the \texttt{affect} of the speaker (i.e., whether the speaker feels positive or negative about the price paid), together with the actual \texttt{price} paid. A speaker talking about $s$ might have the goal of communicating about their \texttt{affect} and the \texttt{price} paid, or they may want to communicate only their \texttt{affect} or only the \texttt{price} paid. Viewed in terms of QUDs, these goals serve to partition the state space. With an \texttt{affect?} QUD, we partition the state space into two cells, corresponding to positive vs.~negative \texttt{affect}. With a \texttt{price?} QUD, we partition the state space into three cells, corresponding to the three possible \texttt{price}s. With a QUD asking after both \texttt{affect-and-price?}, we partition the state space into the six cells listed in \ref{hyperbole-states}.

With the utterance-choice rule in \eqref{S1-QUD}, $S_1$'s goals are QUD specific: for communication to succeed, $L_0$ must simply arrive at the correct cell of the relevant partition; whatever $L_0$ infers about what happened within that partition is irrelevant to $S_1$'s goals. For example, if $S_1$ has a negative affect and is addressing the \texttt{affect?} QUD, utility depends on whether or not $L_0$ correctly arrives at the negative cell of the \texttt{affect?} partition, regardless of whether $L_0$ infers the full state to be $s_1$, $s_2$, or $s_3$. It is at this point where non-literal language suddenly becomes rational for a speaker. To see how, we have to say more about the prior knowledge speakers and listeners bring to bear on their communication scenarios.

World knowledge tells us that states in which the speaker actually paid \$1,000,000 (i.e., $s_3$ and $s_6$) are extraordinarily improbable, while states in which the coffee cost \$3 (i.e., $s_1$ and $s_4$) are the most likely; this knowledge gets reflected in $L_1$'s prior over world states, $P(s)$. We also have prior knowledge about how probable a speaker is to have a positive vs.~negative affect given a specific price paid for a cup of coffee: as prices increase, the probability of negative affect increases along with it. Now, return to a speaker addressing the \texttt{affect?} QUD. If $S_1$'s primary objective is to communicate negative affect to $L_0$, and, crucially, if the available utterances only concern price (i.e., ``I paid \$3/\$7/\$1,000,000''), then the rational thing for $S_1$ to say is that the coffee cost ``\$1,000,000''; given the extremely strong association between \$1,000,000 coffee and negative affect, $L_0$ is all but certain to arrive at the correct answer to the QUD: the speaker holds negative affect.

We have identified how non-literal utterances can be useful to a speaker, but we have yet to capture the non-literal aspect of their meaning. It is at the level of $L_1$ where utterances suddenly become non-literal. Hearing that the speaker paid ``\$1,000,000'', $L_1$ uses the informative priors on prices and their associated affects to arrive at a plausible interpretation. Given that coffee is sure to cost less than \$1,000,000, $L_1$ considers the possible QUDs the speaker might be addressing when selecting $u$. The \texttt{price?} and \texttt{price-and-affect?} QUDs are highly unlikely, given that \$1,000,000 is a highly unlikely price to pay for a cup of coffee. That leaves the \texttt{affect?} QUD. So, $L_1$ considers the possible full states that would have led $S_1$ to utter ``\$1,000,000'' in response to the \texttt{affect?} QUD. Via this counterfactual reasoning, $L_1$ arrives at a reasonable interpretation: $S_1$ is likely to have paid more than usual, so $s_1$ and $s_4$ are downweighted, but we're still on earth, so $s_3$ and $s_6$ remain highly unlikely. Of the remaining states, $S_1$ is most likely to have wanted to communicate negate affect with ``\$1,000,000'', so $L_1$ concludes that $s_5$ is more probably than $s_2$. In other words, $L_1$ hears that the speaker paid ``\$1,000,000'' for a cup of coffee and concludes 1) that the speaker is likely address the \texttt{affect?} QUD, 2) that the speaker is likely trying to communicate negative affect, and 3) that the price is higher than usual but still within the realm of possibility. This last conclusion---a hyperbolic interpretation of the price---is where non-literal meaning enters.

\cite{kaogoodman2015} use an extension of the hyperbole model to capture ironic interpretations of weather descriptions (e.g., ``It's terrible out'' to describe a beautiful sunny day). In addition to inferring speaker affect, the irony model incorporates an additional dimension of meaning that a speaker might want to communicate: their arousal about the weather (i.e., whether they feel strongly about it). The non-literal interpretation mechanism remains the same: if a speaker describes the weather as ``terrible'' but we know that terrible weather is extremely unlikely (the authors situate their imagined conversation in California, where the weather is rarely objectionable), then a listener will conclude that the unlikely literal state (i.e., terrible weather) does not hold, but rather that the speaker must be addressing some other dimension of meaning, namely the strong arousal they feel in their positive affect toward the normal state of affairs: good weather. Similarly in \cite{kaoetal2014metaphor}, who model the interpretation of metaphors like ``John is a whale''. Rather than inferring that John is an aquatic mammal (a highly unlikely state of affairs), the listener concludes that the speaker chose the utterance to communicate about some other dimension of meaning, for example John's physical stature, which has whale-like properties.

\gcs{also discuss the QUD manipulation in the quantifier-scope model?}


\subsection{Context/Prior inference}

We now have technology to model vagueness and ambiguity on the basis of an underspecified semantics, as well as shifting goals and the nonliteral language that addresses them on the basis of uncertainty around the QUD. So far, all of this reasoning has happened with respect to a fixed context, or common ground. What happens when a listener is unsure of the context a speaker has used in the production of their utterance? For example, if you hear that some person is tall, what is the comparison class the speaker had in mind? Depending on the comparison class, the height you ultimately infer for the person  could vary dramatically: tall for a basketball player means something quite different from tall for a jockey.

$$P_{S_1}(u|s, \textbf{x}) \propto \textrm{exp}(\alpha (\textrm{log}P_{L_0}(s|\sem{u}, \textbf{x}) - C(u)))$$

\gcs{MH, could you summarize the comparison class model?}

\gcs{maybe this is also the place to discuss the generics model (in terms of a fancy prior)? or should we mention generics in the context of meaning inference?}

\subsection{Epistemic inference}

- scalar implicature model

- plural predication model


\subsection{Complex utility/utility inference}

- politeness


\section{Modeling practicalities} \label{practicalities}

Free parameters (optimality, cost, alternatives)

World priors: Both how to model and how to measure

QUD

Linking functions


\section{Extensions/limitations} \label{limitations}

Incrementality / processing

Compositionality (CCG chapter in dippl?)

NLP (Dan Klein)

Individual variability (cite Franke \& Degen)

``Just'' a computational level theory. Doesn't bear on mechanisms. But, potentially a hook via resource rational analysis (Leider, Griffiths, ...)

Hand-coded: Don't have a theory of alternatives, state priors

\section{Summary and outlook} \label{summary}

XXX


\bibliography{problang}

%\begin{addresses} %%% uncomment to de-anonymize
%	\begin{address}
%		author1
%		\email{email1}
%	\end{address}
%	\begin{address}
%		author2
%		\email{email2}
%	\end{address}
%	\begin{address}
%		author3
%		\email{email3}
%	\end{address}
%\end{addresses}



\end{document}
